\section{Dataset Audit}
\label{sec:dataset-audit}

% Every number in this section traces to a committed run directory. Sources, by
% subsection:
%   3.1  runs/20260802_class_date_provenance/{class_split_counts.csv,
%        never_evaluated_classes.csv, summary.md}; data/master_results.csv
%   3.2  runs/20260802_class_date_provenance/date_split_summary.csv
%   3.3  runs/20260809_split_ratio_by_group_03/summary.md
%   3.4  runs/20260802_v1_provenance/summary.md;
%        runs/20260801_csv_coverage/split_agreement.csv
%   3.5  runs/20260809_split_v1_vs_v2_by_group_03/summary.md
%   3.6  runs/20260801_device_split/summary.md;
%        runs/20260809_split_v1_vs_v2_by_group_03/summary.md
%   3.7  runs/20260804_duplicate_contamination/summary.md;
%        runs/20260805_figure_verification_sheet_02/summary.md;
%        runs/20260802_burst_clusters_02/gap_histogram.csv
%
% THREE PLACES WHERE THE DRAFTED BULLET AND THE COMMITTED SOURCE DIVERGED. The
% source was followed in each case; see the report accompanying this draft.
%   3.3  the bullet said "within about one point" of 70/20/10. The run
%        classifies on IMAGES (within 1.5 images of a uniform draw) and states
%        explicitly that the percentage-point measure is the wrong one, because
%        it punishes a 54-image session for rounding it cannot avoid. At 1.0 pp
%        only THREE groups qualify, not four. Written on the images measure.
%   3.4  the bullet said T4 is decisive. The run designates T2 as the deciding
%        test and says so in bold: T1 can pass by arithmetic coincidence, so T4
%        rests on T2 having already established identity. Written that way.
%   3.5  the bullet said all eight v1 groups were near-nominal. The run scores
%        SEVEN of eight; the iPhone group misses the tolerance by 1.8 images
%        (0.95 pp) and the run flags it as near-nominal in substance but not by
%        the measure. Written as seven of eight with the borderline stated.
%
% PENDING CITATION -- 3.6. The dataset paper's stated rationale for using
% several phones ("to ensure diverse imaging qualities") cannot be cited:
% references.bib holds only yolov12paper and yolov13paper, with no entry for
% the ElectroCom61 dataset paper itself. The paragraph is written from the
% committed device measurement instead. Add the entry, then restore the
% rationale as a cited clause.

\subsection{Class coverage under the published split}
\label{sec:class-coverage}

The ten most-annotated classes in ElectroCom61, shown at the top of
Figure~\ref{fig:class-instance-counts}, hold at most two test instances each, and
eight of the ten hold none; the eleventh, IC-Base-14-Pin, holds 22.
The transition between the tenth and eleventh ranks is therefore a step rather
than a gradient. This is what distinguishes the pattern from class imbalance,
which is the reading the ordering first invites: an imbalanced dataset produces a
smooth decline in test representation as class frequency falls, and cannot
produce a discontinuity between adjacent ranks in which the more frequent side is
the empty one.

Fifteen of the 61 classes hold zero instances in both the validation and the test
partition, and 22 hold fewer than five test instances. A class with no test
instances cannot contribute to a test-set average precision at all, and a class
with one or two contributes a value estimated from a sample too small to
distinguish a detector from a coin.

Those fifteen classes, listed in Table~\ref{tab:unevaluable-classes}, are not
marginal. They carry 3{,}496 annotation instances between them, 27.0\% of the
12{,}937 in the dataset, and each holds between 153 and 293 training instances.
More than a quarter of the annotation effort in ElectroCom61 therefore supports
classes on which no published result has ever been measured.

The consequence is that every mAP figure reported on this dataset averages over
roughly 46 classes rather than 61, and the two averages are not comparable
quantities. The framework reports the count directly: under the published split
it evaluates 45 classes on validation and 46 on test, across all five detectors
trained here. A score computed over a reduced set of classes is not wrong, but it
answers a narrower question than its name suggests, which is why the count is
reported alongside every figure in Section~\ref{sec:results}.

\begin{figure}[t]
  \centering
  \includegraphics[width=\columnwidth]{f1_class_instance_counts.pdf}
  \caption{Per-class annotation counts, published split. 15 of 61 classes have
  zero valid+test instances (filled marker, bold label) and cannot be evaluated
  at all; 22 in total have fewer than five test instances --- the 15 above plus
  seven more, including IC-Chip, the most annotated class in the dataset.}
  \label{fig:class-instance-counts}
\end{figure}

% Generated by scripts/make_tables.py; carries its own \caption and \label.
\begin{table}[t]
  \centering
  \caption{The 15 classes with zero instances in both validation and test under the published split. Every instance of each sits in training, so none can be evaluated at all. Together they hold 3,496 of the dataset's 12,937 annotations, 27.0\%.}
  \label{tab:unevaluable-classes}
  \begin{tabular}{lrrr}
    \toprule
    Class & Train inst. & Train images & Total inst. \\
    \midrule
    Inductor & 293 & 188 & 293 \\
    LED-Light & 293 & 184 & 293 \\
    MLC-Capacitor & 292 & 81 & 292 \\
    OLED-Display & 279 & 189 & 279 \\
    Push-Switch & 277 & 177 & 277 \\
    Buzzer & 272 & 206 & 272 \\
    High-Voltage-Ceramic-Capacitor & 257 & 200 & 257 \\
    Bluetooth-Module & 249 & 179 & 249 \\
    Sonar-Sensor & 235 & 190 & 235 \\
    Gas-Sensor & 221 & 180 & 221 \\
    LCD-Display & 177 & 177 & 177 \\
    Motor-Driver & 177 & 175 & 177 \\
    LDR-Sensor & 161 & 79 & 161 \\
    9-Volt-Battery & 160 & 159 & 160 \\
    RFID-Scanner & 153 & 153 & 153 \\
    \midrule
    Total (15 classes) & 3,496 & 2,517 & 3,496 \\
    \bottomrule
  \end{tabular}
\end{table}


\subsection{The mechanism}
\label{sec:coverage-mechanism}

Three capture groups sit entirely within the training partition, as panel~(a) of
Figure~\ref{fig:capture-group-composition} shows: 19 February 2024, with 100
images; 20 February 2024, with 486; and the untimestamped iPhone group, with 189. Together they hold 775 images, 52.4\% of the 1{,}478 in the
training set, and none of the three contributes a single image to validation or
test.

All fifteen unevaluable classes occur only within those three groups. The
coverage gap in Section~\ref{sec:class-coverage} is therefore not a property of
the classes but of the sessions they were photographed in: a class that appears
only in a session allocated wholly to training cannot reach an evaluation
partition, however many instances of it the dataset contains. This is also what
makes the gap repairable, because it can be closed by reallocating sessions
rather than by collecting more data --- the construction taken up in
Section~\ref{sec:corrected-split}.

\begin{figure}[t]
  \centering
  \includegraphics[width=\columnwidth]{f2_capture_group_composition.pdf}
  \caption{Capture-group composition of the published split. (a) v2 image counts
  per session: 3 of the 9 groups reach neither valid nor test (bold label, filled
  marker) and together hold 775 of the 1478 training images, 52.4\%; all 15
  classes that cannot be evaluated occur only within them. (b) the same sessions
  before and after the v2 re-split, each normalised to its own total, dashed
  lines where a 70/20/10 split would fall.}
  \label{fig:capture-group-composition}
\end{figure}

\subsection{Allocation is not uniform across sessions}
\label{sec:allocation-uniformity}

Measured session by session, the split takes three distinct shapes. Four groups
--- 6 March, 9 March, 17 April and 18 November 2024 --- fall within 1.5 images of
what a uniform 70/20/10 draw over that group would have produced, which is as
close as integer arithmetic permits. Three contain no validation or test images
at all. The remaining two are strongly skewed in opposite directions: 28 February
2024 holds more validation images than training images, at 35.9/55.2/8.9, and 3
March 2024 holds more test images than training images, at 29.3/20.2/50.5. The
tolerance is expressed in images rather than percentage points because a
54-image session moves 1.9 points per file, so a points-based tolerance would
score it as skewed for rounding it cannot avoid.

Aggregated over all 2{,}121 images the split is 69.7/20.7/9.7, which reads as a
textbook 70/20/10. That aggregate is an accident of offsetting deviations.
Relative to a uniform draw, the three train-only groups withhold 155.0
validation images and 77.5 test images between them; the 28 February group
carries a validation surplus of 169.6 and the 3 March group a test surplus of
76.2. The two nearly cancel, leaving the aggregate 13.8 validation images above
and 7.1 test images below nominal --- under one percentage point in each case.
This is why the problem is invisible at the level at which splits are normally
checked, and it is the reason the per-group measurement was made at all.

\subsection{The metadata file describes v1, not v2}
\label{sec:metadata-provenance}

The metadata CSV shipped inside the v2 archive was tested against an actual v1
download under four provenance tests, and passed all four. The deciding test is
the identity test: the set of image stems in v1 and the set of keys in the CSV
are equal, with zero entries on either side unmatched. Identity is what settles
the question, because the count test can pass by arithmetic coincidence --- two
collections of 2{,}071 files need not be the same files. With identity
established, the split-level test is decisive in turn: the CSV's
\texttt{DATA\_TYPE} column against the folder each image occupies under v1 gives
a perfectly diagonal contingency table, 1{,}454 train, 412 validation and 205
test, with zero disagreements across all 2{,}071 rows.

An independent line of evidence corroborates this. Fifty images in v2 carry no
CSV row at all, and those fifty are split 35/10/5 across train, validation and
test. That is exactly the composition of the 18 November 2024 capture session,
which is the one session absent from v1 entirely. The images the CSV does not
describe are precisely the images that did not exist when it was written.

The consequence is a rule for reading the dataset rather than a defect in it. The
folders on disk are authoritative for the v2 split and the CSV is not: the same
contingency table built against v2 carries 600 off-diagonal rows. But the CSV
remains a faithful record of the v1 assignments, and it is the only surviving
record of them, which is what makes the comparison in the next subsection
possible.

\subsection{What the v2 re-split changed}
\label{sec:v2-resplit}

Under v1, seven of the eight capture groups the CSV covers were near-nominal by
the same 1.5-image measure used above. The eighth, the untimestamped iPhone
group, misses that tolerance by 1.8 images on a 189-image group --- 0.95
percentage points --- so the departure from a uniform draw is a rounding-scale
one in every group present.

The v2 re-split, drawn in panel~(b) of
Figure~\ref{fig:capture-group-composition}, left three of those groups untouched,
with zero images moved in 6 March, 9 March and 17 April, and changed the other
five. In the three groups that
became train-only, validation and test images were folded into training: 20
validation and 10 test images for 19 February, 97 and 48 for 20 February, and 36
and 20 for the iPhone group. The two skewed groups absorbed the resulting
deficits, with 28 February taking the validation shortfall through 189 images
moved from training to validation, and 3 March taking the test shortfall through
77 moved from training to test.

Across the 2{,}071 images the CSV covers, 600 changed assignment between v1 and
v2 --- 29\% of them. These are the same 600 images that produce the off-diagonal
rows in the \texttt{DATA\_TYPE} comparison of
Section~\ref{sec:metadata-provenance}; the disagreement between the CSV and the
folders is a record of the re-split, not an inconsistency within either.

One further observation constrains what can be inferred from the above. The 18
November 2024 session, which was added in v2 and has no v1 assignment to change,
was split 35/10/5 of 50 --- exactly 70/20/10, with zero deviation in any of the
three cells. Whatever produced the v2 split allocated a new session at the
nominal ratio. No mechanism is claimed here for the changes to the pre-existing
assignments; the observation is recorded because it bounds the explanations that
remain available.

\subsection{One camera is never evaluated}
\label{sec:camera-coverage}

Under v1, the 189 untimestamped images, which the metadata attributes to an
iPhone, were split 133\slash 36\slash 20 across train, validation and test. Under
v2 they are split 189\slash 0\slash 0, so every image from that device now sits
in the training partition.

The dataset records four distinct device names across its 2{,}121 images, and the
iPhone is the only one confined to a single partition. Whatever variation in
imaging characteristics that device contributes --- sensor, optics, and the
processing pipeline applied before the image was ever exported --- is therefore
present in what the detectors learn from and absent from what they are measured
on. A model that generalises poorly to that device would not register as
degraded under the published split, because no test image would expose it.

% TODO: the dataset paper states that several phones were used to ensure diverse
% imaging qualities, which is the sentence that makes this subsection's
% implication sharp. It cannot be cited until references.bib carries an entry
% for the dataset paper. Restore as a cited clause once it does.

\subsection{Duplicate contamination: a negative result}
\label{sec:duplicate-contamination}

The published split carries no test--train near-duplicate pairs at any threshold
tested, under either of the two scorings applied, with low-information pairs
excluded. This is reported as a finding rather than as an absence of one. The
split has a documented coverage defect, established above, and the natural
suspicion is that a split allocated in a way that stranded three whole sessions
in training might also have scattered near-identical frames across the boundary.
It did not, and a corrected split must therefore preserve this property rather
than merely avoid making it worse.

The margin behind that result is thin and is disclosed here rather than left to
be discovered. The closest test--train pair scores 0.0515 against the loosest
threshold applied, 0.05, clearing it by 3\%. A threshold chosen slightly more
conservatively would have returned a non-zero count, so the result is
threshold-sensitive and should be quoted with the threshold attached.

No threshold can be derived empirically from the data itself. The distribution of
gaps between consecutive captures is unimodal: measured per second of bin width,
it rises monotonically to a peak in the 3--5 second bin and decays monotonically
thereafter, on both of the two independent device keyings tested. There is no
valley between a burst mode and a between-burst mode at which a cut could be
placed, because there is only one mode. Any grouping threshold is consequently a
construction rather than a measurement, which is why the one adopted in
Section~\ref{sec:corrected-split} is justified by the properties of the split it
produces rather than by the shape of this distribution.

Two scope limits travel with every figure above. Only pairs sharing an identical
class inventory are compared, so a genuine duplicate in which one component is
occluded in one frame is never scored; every count is a floor rather than an
estimate. And 95 low-information cross-split pairs, those carrying two annotated
boxes or fewer, are excluded from the ranking, because a match between box
centres is cheap to achieve by chance when there are few boxes to match.
